\documentclass[12pt,a4paper]{article}
\usepackage{amsmath, amsfonts, amssymb}
\usepackage{geometry}
\usepackage{multirow}
\usepackage{fancyhdr}
\usepackage{array}

% Setting margin presisi
\geometry{a4paper, top=2.5cm, bottom=2.5cm, left=2cm, right=2cm, headheight=15pt}

% Bikin custom fancy header buat halaman soal aja (mulai hal 3)
\fancypagestyle{halamansoal}{
    \fancyhf{}
    \fancyhead[L]{Nama/NRP: \hrulefill}
    \renewcommand{\headrulewidth}{0pt}
    \cfoot{\thepage}
}

\begin{document}

% ==================== HALAMAN 1 (COVER) ====================
\thispagestyle{empty}

\begin{center}
  {\Large \textbf{OLIMPIADE SAINS ITS (OSITS) 2026}}\\[0.2cm]
  {\Large \textbf{BABAK PENYISIHAN}}
\end{center}

\vfill

\begin{center}
  {\Large \textbf{BIDANG MATEMATIKA}}\\[0.2cm]
  {\Large \textbf{Sabtu, 28 Februari 2026}}
\end{center}

\vfill

\begin{center}
  {\Large \textbf{Waktu : 180 menit (3 jam)}}
\end{center}

\vfill
\vfill

\noindent
\begin{tabular}{|p{2.2cm}|p{1.4cm}|p{1.4cm}|p{1.4cm}|p{1.4cm}|p{1.4cm}|p{3.5cm}|}
  \hline
  \rule[-0.2cm]{0pt}{0.7cm}Nama :   & \multicolumn{6}{l|}{}                                                             \\
  \hline
  \rule[-0.2cm]{0pt}{0.7cm}NRP :    & \multicolumn{6}{l|}{}                                                             \\
  \hline
  \multirow{4}{*}{\vspace{1cm}SKOR} & \rule{0pt}{0.4cm}1.   & 2. & 3. & 4. & 5.  & \centering\arraybackslash TOTAL SKOR \\ \cline{7-7}
                                    & \rule{0pt}{0.7cm}     &    &    &    &     & \multirow{3}{*}{}                    \\ \cline{2-6}
                                    & \rule{0pt}{0.4cm}6.   & 7. & 8. & 9. & 10. &                                      \\
                                    & \rule{0pt}{0.8cm}     &    &    &    &     &                                      \\
  \hline
\end{tabular}

\newpage

% ==================== HALAMAN 2 (PETUNJUK) ====================
\thispagestyle{empty}
\begin{center}
  \textbf{PETUNJUK PENGERJAAN}\\
  \textbf{OSITS BIDANG MATEMATIKA}
\end{center}

\begin{enumerate}
  \item OSITS bidang Matematika terdiri 10 soal uraian dengan maksimal skor untuk setiap soalnya adalah 10 poin. Bidang yang diujikan meliputi:
        \begin{itemize}
          \item Analisis Real
          \item Aljabar Linier
          \item Struktur Aljabar
          \item Analisis Kompleks
          \item Kombinatorika
        \end{itemize}
  \item Waktu pengerjaan selama 180 menit dan tidak ada penambahan waktu bagi yang datang terlambat.
  \item Tuliskan nama dan NRP di sampul bagian depan dan setiap halaman soal (bagian atas).
  \item Tuliskan jawaban Saudara lengkap dengan argumentasi dan penjelesan.
  \item Tuliskan jawaban Saudara pada kotak jawaban yang tersedia di setiap soalnya. Jika tidak mencukupi diperbolehkan menggunakan halaman bagian belakang.
  \item Wajib menggunakan bolpoin untuk menuliskan jawaban.
  \item Dilarang menggunakan kalkulator atau alat bantu hitung lainnya selama pengerjaan berlangsung.
  \item Dilarang berbuat curang seperti bekerja sama, mencontek, ataupun jenis pelanggaran lainnya.
  \item Keputusan juri tidak dapat diganggu gugat.
\end{enumerate}

\newpage

% MENGAKTIFKAN HEADER NAMA/NRP UNTUK HALAMAN SOAL
\pagestyle{halamansoal}

% Buka block enumerate utama untuk soal 1-10
% ==================== HALAMAN 3 (SOAL 1) ====================
\begin{enumerate}
  \item \textbf{(Analisis Real)}\\
        Buktikan bahwa untuk setiap bilangan real $a>0$ berlaku
        \[ \ln(1+a) \le \frac{a}{\sqrt{1+a}} \]


        \textbf{Jawaban:}
        \par\vspace{0.1cm}\noindent
        \fbox{%
          \begin{minipage}[t][\dimexpr\textheight-\pagetotal-2.5cm\relax]{\dimexpr\linewidth-2\fboxsep-2\fboxrule\relax}
            \mbox{}
          \end{minipage}%
        }

        \newpage

        % ==================== HALAMAN 4 (SOAL 2) ====================
  \item \textbf{(Aljabar Linier)}\\
        Diberikan bilangan real $a, b, c, d$ sehingga $c \ne 0$ dan $ad-bc=1$. Buktikan bahwa terdapat bilangan real $k, u,$ dan $v$ sehingga
        \[
          \begin{pmatrix} a & b \\ c & d \end{pmatrix} =
          \begin{pmatrix} k & -u \\ 0 & k \end{pmatrix}
          \begin{pmatrix} k & 0 \\ c & k \end{pmatrix}
          \begin{pmatrix} k & -v \\ 0 & k \end{pmatrix}.
        \]


        \textbf{Jawaban:}
        \par\vspace{0.1cm}\noindent
        \fbox{%
          \begin{minipage}[t][\dimexpr\textheight-\pagetotal-2.5cm\relax]{\dimexpr\linewidth-2\fboxsep-2\fboxrule\relax}
            \mbox{}
          \end{minipage}%
        }

        \newpage

        % ==================== HALAMAN 5 (SOAL 3) ====================
  \item \textbf{(Struktur Aljabar)}\\
        Misalkan $G$ grup komutatif dan $g \in G$ dengan $|g|=2$. Jika $H$ subgrup dari $G$, buktikan bahwa $H \cup gH$ juga subgrup dari $G$. Apakah $H \cup gH$ subgrup normal? Jelaskan.\\


        \textbf{Jawaban:}
        \par\vspace{0.1cm}\noindent
        \fbox{%
          \begin{minipage}[t][\dimexpr\textheight-\pagetotal-2.5cm\relax]{\dimexpr\linewidth-2\fboxsep-2\fboxrule\relax}
            \mbox{}
          \end{minipage}%
        }

        \newpage

        % ==================== HALAMAN 6 (SOAL 4) ====================
  \item \textbf{(Analisis Kompleks)}\\
        Diberikan bilangan kompleks $z = \cos \theta + i \sin \theta$ dengan $z \ne 1$. Tunjukkan bahwa
        \[ \text{Re}\left(\frac{1}{1-z}\right) = \frac{1}{2} \]


        \textbf{Jawaban:}
        \par\vspace{0.1cm}\noindent
        \fbox{%
          \begin{minipage}[t][\dimexpr\textheight-\pagetotal-2.5cm\relax]{\dimexpr\linewidth-2\fboxsep-2\fboxrule\relax}
            \mbox{}
          \end{minipage}%
        }

        \newpage

        % ==================== HALAMAN 7 (SOAL 5) ====================
  \item \textbf{(Kombinatorika)}\\
        Misalkan $M$ adalah himpunan semua bilangan 8-digit yang digit-digitnya adalah 1, 2 atau 3 dan memuat paling sedikit 1 digit 2. Tentukan Banyaknya bilangan $N$ di $M$ sehingga setiap digit 2 di $N$ diapit oleh digit 1 dan 3.\\


        \textbf{Jawaban:}
        \par\vspace{0.1cm}\noindent
        \fbox{%
          \begin{minipage}[t][\dimexpr\textheight-\pagetotal-2.5cm\relax]{\dimexpr\linewidth-2\fboxsep-2\fboxrule\relax}
            \mbox{}
          \end{minipage}%
        }

        \newpage

        % ==================== HALAMAN 8 (SOAL 6) ====================
  \item \textbf{(Analisis Real)}\\
        Misalkan $f:[0,1] \rightarrow \mathbb{R}$ terdiferensialkan dan turunan $f^{\prime}$ bersifat Lipschitz pada $[0,1]$ dengan konstanta 6, yaitu
        \[ |f^{\prime}(x)-f^{\prime}(y)| \le 6|x-y| \quad \text{untuk semua } x, y \in [0,1] \]
        Diketahui pula bahwa
        \[ f(0)=2, \quad f(1)=1 \]
        Buktikan bahwa untuk setiap $x \in [0,1]$ berlaku
        \[ |f(x)-(2-x)| \le 3x(1-x). \]


        \textbf{Jawaban:}
        \par\vspace{0.1cm}\noindent
        \fbox{%
          \begin{minipage}[t][\dimexpr\textheight-\pagetotal-2.5cm\relax]{\dimexpr\linewidth-2\fboxsep-2\fboxrule\relax}
            \mbox{}
          \end{minipage}%
        }

        \newpage

        % ==================== HALAMAN 9 (SOAL 7) ====================
  \item \textbf{(Aljabar Linier)}\\
        Misalkan $T:\mathbb{R}^{2} \rightarrow \mathbb{R}^{2}$ adalah pemetaan linier dan untuk setiap $v \in \mathbb{R}^{2}$ berlaku
        \[ ||Tv|| = ||v||. \]
        Buktikan bahwa $T$ adalah transformasi ortogonal. Dengan kata lain, jika $A$ adalah matriks representasi $T$ pada basis standar ($Tv=Av$), tunjukkan bahwa
        \[ A^{T}A = I. \]


        \textbf{Jawaban:}
        \par\vspace{0.1cm}\noindent
        \fbox{%
          \begin{minipage}[t][\dimexpr\textheight-\pagetotal-2.5cm\relax]{\dimexpr\linewidth-2\fboxsep-2\fboxrule\relax}
            \mbox{}
          \end{minipage}%
        }

        \newpage

        % ==================== HALAMAN 10 (SOAL 8) ====================
  \item \textbf{(Struktur Aljabar)}\\
        Diberikan $b=(1~3~5~4)(2~6~3~5) \in S_{15}$.
        \begin{enumerate}
          \item Dapatkan $b^{2026}$
          \item Tentukan banyaknya $m \in S_{15}$ yang memenuhi $b=m^{5}$.
        \end{enumerate}


        \textbf{Jawaban:}
        \par\vspace{0.1cm}\noindent
        \fbox{%
          \begin{minipage}[t][\dimexpr\textheight-\pagetotal-2.5cm\relax]{\dimexpr\linewidth-2\fboxsep-2\fboxrule\relax}
            \mbox{}
          \end{minipage}%
        }

        \newpage

        % ==================== HALAMAN 11 (SOAL 9) ====================
  \item \textbf{(Analisis Kompleks)}\\
        Diberikan polinomial $f(x)=x^{4}+ax^{3}+bx^{2}+cx+d$ dimana setiap koefisiennya merupakan bilangan real dan akarnya real. Buktikan bahwa jika $|f(i)|=1$ maka $a=b=c=d=0$.\\


        \textbf{Jawaban:}
        \par\vspace{0.1cm}\noindent
        \fbox{%
          \begin{minipage}[t][\dimexpr\textheight-\pagetotal-2.5cm\relax]{\dimexpr\linewidth-2\fboxsep-2\fboxrule\relax}
            \mbox{}
          \end{minipage}%
        }

        \newpage

        % ==================== HALAMAN 12 (SOAL 10) ====================
  \item \textbf{(Kombinatorika)}\\
        Tentukan bentuk sederhana dari
        \[ \sum_{j=0}^{k} 2^{k-j} \binom{k+j}{j} \]


        \textbf{Jawaban:}
        \par\vspace{0.1cm}\noindent
        \fbox{%
          \begin{minipage}[t][\dimexpr\textheight-\pagetotal-2.5cm\relax]{\dimexpr\linewidth-2\fboxsep-2\fboxrule\relax}
            \mbox{}
          \end{minipage}%
        }

\end{enumerate} % Tutup block enumerate utama
\end{document}